% Section 1: Introduction

Overconfidence---the tendency to overestimate the precision of one's information---is a well-documented behavioral bias in financial markets and is empirically associated with excessive trading and episodes of mispricing \cite{Odean1998,BarberOdean2001}. A central modeling challenge is to combine heterogeneous (and potentially biased) beliefs formed from private information with endogenous price formation, while retaining an equilibrium concept that is mathematically well-posed, computable, and connectable to data. This paper develops a mean field game (MFG) framework for a market populated by a continuum of investors who filter an unobserved fundamental from noisy private signals and trade against an endogenous price impacted by mean trading rate (order flow).

The model keeps a clean stock--flow separation. Inventory $x_{i,t}$ (shares) carries over between decision times, trading rate $u_{i,t}$ (shares per unit time) satisfies $x_{i,t+\Delta t}=x_{i,t}+u_{i,t}\Delta t$, and the price equation loads on the mean flow $\bar u_t$ rather than on mean inventory. Overconfidence enters only through \emph{filtering}: agents treat their private signal as more precise than it truly is, which changes posterior uncertainty and therefore trading intensity and disagreement, not through an ad hoc demand multiplier.

Early behavioral models introduced stylized agent types (noise traders, momentum traders) to demonstrate how bounded rationality can drive prices away from fundamentals \cite{DeLong1990,Hong1999,DanielHirshleiferSubrahmanyam1998,GervaisOdean2001}, but typically rely on a small number of types rather than a full distribution of beliefs. Mean field games have been applied to financial markets in several contexts: Gu\'eant, Lasry, and Lions \cite{GueantLasryLions2010} developed MFG frameworks for execution and portfolio problems; Carmona, Fouque, and Sun \cite{CarmonaFouqueSun2015} studied systemic risk; Lachapelle \emph{et al.} \cite{LachapelleLasryLehalleLions2016} applied MFG ideas to high-frequency trading and price formation; and Cardaliaguet and Lehalle \cite{CardaliaguetLehalle2018} analyzed trade crowding effects.

Our contribution relative to existing MFG finance literature is to provide a bounded-rational equilibrium concept that is both economically interpretable and mathematically checkable in a stock--flow microstructure setting. In the dynamic baseline, we study a discrete-time \emph{myopic mean-field equilibrium}: given a conjectured common-noise adapted mean trading rate $(\bar u_t)$, each agent applies a one-step \CARA best response under conditionally Gaussian price increments (using their filtered belief), and mean-field consistency closes the dynamics by requiring $\bar u_t$ to equal the conditional population mean of these best responses. Because the best response is linear in $(\hat v_{i,t}-p_t)$ and in $\bar u_t$, the closure reduces each time step to a scalar stock--flow fixed point. A per-step contraction condition guarantees uniqueness and well-posedness of the induced discrete-time dynamics and particle approximation (Theorem~\ref{thm:contraction}).

The simulations distinguish \emph{level amplification} (large mispricing/volatility levels driven by weak anchoring and high noise trading) from \emph{overconfidence effects} (incremental differences when increasing $k$). Across a targeted joint grid over anchoring, impact, and noise trading, we find a robust mechanism separation in the baseline anchored-impact regime: microstructure parameters govern mispricing levels, while overconfidence strongly increases disagreement and trading intensity with economically small effects on price-level mispricing (Subsection~\ref{subsec:joint_grid}). To make the model empirically usable, Section~\ref{sec:param_identification} provides an identification argument linking key parameters to observable price and volume moments, and a moment-match vignette illustrates the calibration bridge on a liquid ETF.

The paper's dynamic analysis has a deliberate two-model structure. (A) The object we compute and simulate is the discrete-time stock--flow myopic equilibrium described above. (B) As a theoretical benchmark, Subsection~\ref{subsec:ct_cara_lqg} derives an exact finite-horizon continuous-time \CARA/\LQG\ best response under partial information; this analytic extension clarifies how the full intertemporal solution adds a hedging correction to the myopic rule, and we compare the myopic and intertemporal policies in Subsection~\ref{subsec:extensions}.

The main contributions of this paper are:
\begin{itemize}
    \item We formulate a mean field game with overconfident Bayesian filtering and an explicit stochastic price equation with endogenous flow impact, under a stock--flow timing that separates inventory and trading rate (Subsubsection~\ref{subsubsec:timing_stock_flow}).
    \item We provide a clean static benchmark with a closed-form linear price \eqref{eq:static_price} and explicit comparative statics in the overconfidence parameter $k$.
    \item In the dynamic setting, we derive a checkable myopic \CARA inventory rule under conditionally Gaussian price increments (with holding cost) and obtain a per-step stock--flow mean-field closure; an explicit contraction condition guarantees uniqueness and discrete-time well-posedness (Theorem~\ref{thm:contraction}).
    \item We report particle-based numerical experiments (large-$N$ simulations) across bull/bear regimes and microstructure regimes, including mechanism diagnostics (belief dispersion and perceived uncertainty), validation checks, and extensions (Section~\ref{sec:numerical}).
    \item We provide a minimal identification/calibration bridge: Section~\ref{sec:param_identification} links key parameters to observable price and volume moments and illustrates the mapping with a moment-match vignette.
\end{itemize}

The remainder of this paper is organized as follows. Section~\ref{sec:model} presents the model, including the static benchmark and the dynamic price/belief dynamics. Section~\ref{sec:meanfield} discusses the mean-field object in the presence of common noise and the fixed-point notion of (myopic) equilibrium. Section~\ref{sec:results} states the main analytic results, including the myopic rule and the continuous-time \CARA/\LQG\ benchmark. Section~\ref{sec:proofs} collects proofs. Section~\ref{sec:numerical} presents numerical experiments, the empirical bridge, and the joint-grid analysis. Section~\ref{sec:conclusion} concludes.
